\documentclass{article}
\begin{document}

\begin{flushright}
Jiancheng Liu\\
College of Chemical and Biological Engineering\\
Zhejiang University, Hangzhou 310058, China\\
Email: 3240101908@zju.edu.cn
\end{flushright}

\vspace{1cm}

\noindent Dear Editor,

I am pleased to submit our manuscript entitled \textbf{``Multi-Mechanism Pre-Training of Physics-Informed Neural Networks for Fouling Resistance Reconstruction from Sparse Measurements''} for consideration in \textit{Thermal Science and Engineering Progress}.

\vspace{0.5cm}

\noindent \textbf{Summary and contribution}

This paper presents an exploratory computational study on using multi-mechanism physics-informed pre-training to reconstruct heat exchanger fouling resistance curves from sparse measurement data. We propose a framework where a neural network is pre-trained on diverse synthetic data generated from five established fouling ODE models, then fine-tuned on as few as 5--30 real data points. The framework is validated on 18 fouling curves digitized from three peer-reviewed studies. The key finding is that \textit{mechanism diversity} during pre-training prevents the catastrophic over-specialization observed when pre-training on a single ODE model, providing performance gains that cannot be replicated by simply increasing data volume.

We report results transparently, including competitive baselines (smoothing spline, Gaussian Process regression) that match or exceed the PINN on pure interpolation tasks. The framework's principal value lies in providing physics-grounded representations that connect sparse reconstruction to downstream cleaning decisions---a workflow that classical interpolation methods cannot offer. All code, data, and pre-trained model weights are publicly available at \url{https://github.com/wy2627962897/fouling-pinn-pretrain}.

\vspace{0.5cm}

\noindent \textbf{Why TSEP?}

This manuscript aligns with TSEP's scope in thermal process modelling, simulation, and optimization. It addresses a practical engineering problem---fouling resistance estimation under data scarcity---using a computational approach that bridges established physical models with modern machine learning techniques. The topic is well suited to TSEP's readership in thermal engineering applications and industrial energy efficiency.

\vspace{0.5cm}

\noindent \textbf{Declarations}
\begin{itemize}
  \item This manuscript has not been published previously and is not under consideration elsewhere.
  \item All authors have approved the manuscript and agree with its submission to TSEP.
  \item The author declares no competing financial interests or personal relationships that could have influenced this work.
  \item This research did not receive any specific grant from funding agencies in the public, commercial, or not-for-profit sectors.
\end{itemize}

\vspace{0.5cm}

\noindent Thank you for your consideration.

\vspace{0.5cm}

\noindent Sincerely,\\
Jiancheng Liu

\end{document}
